% ******************************************************** %
%  TP FINAL -- MICROKERNEL RISC-V                          %
%  Portación del TP3 "Crystal Campus" de x86 a RISC-V rv64 %
% ******************************************************** %

\documentclass[a4paper]{article}
\usepackage[spanish]{babel}
\usepackage[utf8]{inputenc}
\usepackage{charter}   % tipografía
\usepackage{graphicx}
\usepackage{wrapfig}
\usepackage{fancyvrb}
\usepackage{paralist}
\usepackage{amsmath, amsthm, amssymb}
\usepackage{multicol}
\usepackage{underscore}
\usepackage{caratula}
\usepackage{url}
\usepackage{booktabs}
\usepackage{longtable}

% ********************************************************* %
% ~~~~~~~~              Code snippets             ~~~~~~~~~ %
% ********************************************************* %
\usepackage{color}
\usepackage{fancybox}
\definecolor{litegrey}{gray}{0.94}
\newenvironment{codesnippet}{%
    \begin{Sbox}\begin{minipage}{\textwidth}\sffamily\small}%
    {\end{minipage}\end{Sbox}%
        \begin{center}%
        \vspace{-0.4cm}\colorbox{litegrey}{\TheSbox}\end{center}\vspace{0.3cm}}

% ********************************************************* %
% ~~~~~~~~         Formato de las páginas         ~~~~~~~~~ %
% ********************************************************* %
\usepackage{fancyhdr}
\pagestyle{fancy}
\fancyhf{}
\fancyhead[LO]{\small Organización del Computador II}
\fancyhead[RO]{\small Microkernel RISC-V}
\fancyfoot[LO]{\small{Joaquín Romera}}
\fancyfoot[RO]{\thepage}
\renewcommand{\headrulewidth}{0.5pt}
\renewcommand{\footrulewidth}{0.5pt}
\setlength{\hoffset}{-0.8in}
\setlength{\textwidth}{16cm}
\setlength{\headsep}{0.5cm}
\setlength{\textheight}{25cm}
\setlength{\voffset}{-0.7in}
\setlength{\headwidth}{\textwidth}
\setlength{\headheight}{13.1pt}
\renewcommand{\baselinestretch}{1.1}

% Mayor "estiramiento" de párrafos para \texttt{} largos sin hifenar.
\emergencystretch=3em
\tolerance=800
\hbadness=10000

% ******************************************************** %

\begin{document}

% Desactivamos shorthands de babel-spanish (<<, >>, ~) para evitar conversión
% a guillemets dentro de snippets y \texttt{}.
\shorthandoff{<>~}

\thispagestyle{empty}
\materia{Organización del Computador II}
\submateria{TP Final}
\titulo{Microkernel RISC-V}
\subtitulo{Adaptación del TP3 \emph{Crystal Campus} de x86 (IA-32) a RISC-V (rv64)}
\fecha{\today}
\integrante{Joaquín Romera}{183/16}{joakromera@gmail.com}

\maketitle
\newpage

\thispagestyle{empty}
\vfill
\vspace{3cm}

\begin{abstract}
Este informe documenta la adaptación del TP3 \emph{Crystal Campus} de Organización del Computador II, originalmente
implementado sobre x86 (IA-32) en modo protegido y corrido bajo Bochs, a la arquitectura RISC-V de 64 bits
(\texttt{rv64imac}) sobre la máquina \texttt{virt} de QEMU. El objetivo es preservar la mecánica de juego y el esqueleto
de ejercicios del TP original, reemplazando los mecanismos propios de IA-32 (GDT, IDT, TSS, PIC 8259A, PIT 8254, teclado
PS/2, memoria de video VGA en \texttt{0xB8000}, paginación de dos niveles con CR0/CR3) por sus equivalentes nativos de
RISC-V (Sv39, \texttt{stvec}, \texttt{satp}, CLINT/ACLINT, UART 16550, modos M/S/U). El kernel corre sin firmware previo
(QEMU \texttt{-bios none}): arranca en \textbf{M-mode}, se configura él mismo (PMP, delegación de traps, timer) y
desciende a \textbf{S-mode}, donde vive el grueso del kernel; las tareas mineras corren en \textbf{U-mode}. La arquitectura
distintiva del TP (tener el código de cada minero físicamente alojado en la celda del tablero donde se encuentra)
se preserva sobre Sv39, manteniendo la noción de ``tablero como memoria''.
\end{abstract}

\newpage

\tableofcontents

\section{Introducción}

El TP3 de la materia se implementaba originalmente sobre x86 (IA-32) en modo protegido, con boot desde floppy, BIOS y
ejecución bajo el simulador Bochs. Este trabajo final documenta la adaptación de ese mismo proyecto a la arquitectura
RISC-V de 64 bits (\texttt{rv64imac}) sobre la máquina \texttt{virt} de QEMU.

La idea de este trabajo es preservar la mecánica de juego y el esqueleto de ejercicios del TP original,
reemplazando los mecanismos específicos de IA-32 (GDT, IDT, TSS, PIC 8259A, PIT 8254, teclado PS/2, memoria de video VGA
en \texttt{0xB8000}, paginación de dos niveles con CR0/CR3) por sus equivalentes nativos de RISC-V. En particular, la
consigna principal del TP --- tener el código de cada minero \emph{físicamente alojado} en la celda del tablero donde
se encuentra y, al moverse, copiar su imagen de 8~KiB a la nueva celda y remapear \texttt{TASK\_CODE} --- se mantiene
uno a uno, pero ahora implementada sobre Sv39.

El kernel corre sin firmware previo (QEMU \texttt{-bios none}). Esto implica que el propio kernel arranca en
\textbf{modo máquina (M-mode)}, se configura a sí mismo (PMP, delegación de traps, tick del timer) y desciende luego a
\textbf{modo supervisor (S-mode)}, donde vive el grueso del kernel. Las tareas mineras corren en \textbf{modo usuario
(U-mode)}. Los detalles de los CSRs y modos de privilegio siguen~\cite{riscv-priv}.

Este informe está organizado siguiendo la misma numeración de ejercicios del enunciado y del informe x86 original, para
que la correspondencia entre ambos sea uno a uno. En cada sección se indica, primero, qué se hacía en x86 y, a
continuación, qué mecanismo nativo de RISC-V lo reemplaza y dónde está implementado dentro de la carpeta \texttt{riscv/} del
repositorio.

\newpage
\section{Entorno y herramientas}

\subsection{De Bochs a QEMU \texttt{virt}}

El TP original corría sobre Bochs, emulando un IBM-PC compatible con boot desde floppy y BIOS. En RISC-V esto no
aplica; la opción estándar para bare-metal es la máquina \texttt{virt} de QEMU~\cite{qemu-virt}, que expone:

\begin{itemize}
  \item Un conjunto de HARTs rv64 (este port usa uno solo).
  \item RAM a partir de \texttt{0x80000000}.
  \item Un UART 16550~\cite{ns16550} mapeado en \texttt{0x10000000}.
  \item CLINT (Core Local Interruptor) con \texttt{mtime}/\texttt{mtimecmp} en el rango
  \texttt{0x02000000}~\cite{riscv-aclint}.
  \item PLIC en \texttt{0x0C000000} (no se usa en este port; el tick sale del CLINT).
\end{itemize}

Al arrancar con \texttt{-bios none}, QEMU carga el ELF del kernel y salta a \texttt{0x80000000} en M-mode, sin
firmware intermedio. Esto deja al kernel la responsabilidad de todo el boot, incluyendo la programación del timer.
El wiki OSDev~\cite{osdev-riscv} documenta en detalle esta forma de arranque y el layout de MMIO de la máquina
\texttt{virt}.

\subsection{Toolchain y flags}

Se utiliza la toolchain \texttt{riscv64-unknown-elf-*} (GCC bare-metal) y \texttt{qemu-system-riscv64}. Los flags
relevantes del \texttt{Makefile} del port:

\begin{codesnippet}
\begin{verbatim}
 ARCH_OPTS  = -march=rv64imac -mabi=lp64
 COMMONFLAGS = -Wall -O2 -ffreestanding -nostdlib -nostartfiles \
               -fno-stack-protector -fno-exceptions -fno-rtti \
               -fno-threadsafe-statics -fno-unwind-tables \
               -fno-asynchronous-unwind-tables -mcmodel=medany $(ARCH_OPTS)
\end{verbatim}
\end{codesnippet}

El \emph{code model} \texttt{medany} es necesario porque el kernel se ubica en \texttt{0x80000000}, muy fuera del
rango $\pm 2$~GiB del PC que usa \texttt{medlow} para instrucciones PC-relativas extendidas~\cite{riscv-psabi}. El
script del linker (\texttt{riscv/linker.ld}) fija \texttt{.}~$=$~\texttt{0x80000000} y coloca \texttt{\_start} primero,
de manera que el punto de entrada coincide con el reset vector de QEMU.

\newpage
\section{Mapeo general x86 $\to$ RISC-V}
\label{sec:mapeo}

La siguiente tabla resume las equivalencias que se irán desarrollando en las secciones siguientes. Varios elementos del
TP3 no tienen contraparte directa y su función la cumple otro mecanismo; se indica con ``---''.

\begin{center}
\small
\begin{tabular}{p{6.2cm} | p{8.6cm}}
\textbf{TP3 x86} & \textbf{RISC-V (este port)} \\
\hline
GDT + segmentación (code/data, nivel 0/3) & --- (no hay segmentos). Permisos por PTE Sv39 (\texttt{U/R/W/X}). \\
Modo real $\to$ modo protegido + A20 & M-mode $\to$ S-mode vía \texttt{mret}; PMP abierta. \\
CR0.PG + CR3 & CSR \texttt{satp} con modo SV39 + ASID + PPN. \\
IDT + 255 ISR stubs & CSR \texttt{stvec} (único) + despacho en C por \texttt{scause}. \\
Stack separada en cambio de nivel (SS0:ESP0 en TSS) & Stack dedicada de trap (\texttt{trap\_stack\_top}) y salto de
\texttt{gp} al del kernel. \\
\texttt{pushad}, selectores + \texttt{iret} & Guardado manual en \texttt{TrapFrame} (GPRs +
\texttt{sepc}/\texttt{sstatus}/\texttt{scause}/\texttt{stval}) y \texttt{sret}. \\
PIT 8254 + IRQ0 + PIC \texttt{pic\_finish1} & CLINT \texttt{mtimecmp}/\texttt{mtime} en M-mode $\to$ SSIP; \texttt{csrc
sip, SSIP}. \\
Teclado PS/2 IRQ1 + scancodes & UART 16550 por polling; teclas ASCII. \\
VGA texto en \texttt{0xB8000} & UART + secuencias ANSI (80$\times$50 emulado). \\
TSS + cambio de tarea por hardware (\texttt{jmp <sel>:<off>}) & \texttt{TrapFrame} por tarea + \texttt{sret} (cambio por
software). \\
\texttt{int 0x2F} (syscall) & \texttt{ecall}: \texttt{a7}=id, \texttt{a0}=arg/retorno. \\
\texttt{invlpg} / recargar CR3 & \texttt{sfence.vma} (por VA o global, opcionalmente por ASID). \\
\texttt{bochsrc} / Bochs & QEMU \texttt{virt} con \texttt{-bios none}. \\
\end{tabular}
\end{center}

Las referencias a instrucciones y CSRs siguen la convención de~\cite{riscv-unpriv} (ISA no privilegiada)
y~\cite{riscv-priv} (ISA privilegiada, CSRs, Sv39).

\newpage
\section{Ejercicio 1 --- Arranque y paso a modo supervisor}
\subsection{Qué hacía el TP x86}

El TP3 original arrancaba en modo real desde el boot-sector, habilitaba A20, cargaba la GDT con \texttt{lgdt}, prendía PE de CR0,
hacía el \texttt{jmp} lejano a modo protegido, cargaba los segmentos y fijaba el stack en \texttt{0x27000}. Además declaraba
un segmento de pantalla aparte y limpiaba la VGA.

\subsection{Equivalente en RISC-V}

En RISC-V no hay segmentación ni A20. Los permisos de memoria salen de la PTE Sv39 (\texttt{V}, \texttt{R}, \texttt{W},
\texttt{X}, \texttt{U}, \texttt{G}, \texttt{A}, \texttt{D}), y lo más parecido a ``pasar a modo protegido'' es
\textbf{ejecutar \texttt{mret} desde M-mode con \texttt{mstatus.MPP}$=$1 (S-mode) y
\texttt{mepc}$=$\texttt{kernel\_main}}.

La cadena de arranque es:

\begin{enumerate}
  \item QEMU salta a \texttt{\_start} en \texttt{0x80000000} (M-mode).
  \item \texttt{entry.S} inicializa \texttt{gp}, stack del kernel y limpia \texttt{.bss}, y llama a
  \texttt{machine\_init}.
  \item \texttt{machine.S::machine\_init} (una ``mini firmware'' en M-mode, ver sección~\ref{sec:machine}) configura
  PMP, delega traps a S-mode, programa el tick y ejecuta \texttt{mret} $\to$ \texttt{kernel\_main} ya en S-mode.
\end{enumerate}

\subsection{\texttt{riscv/entry.S}}

Fragmento relevante del entrypoint:

\begin{codesnippet}
\begin{verbatim}
 _start:
     .option push
     .option norelax
     la gp, __global_pointer$
     .option pop

     csrw mie, zero
     csrci mstatus, 0x8        /* MIE */
     csrw sie, zero
     csrci sstatus, 0x2        /* SIE */

     la sp, kernel_stack_top

     /* Clear .bss */
     la t0, __bss_start
     la t1, __bss_end
 1:
     bgeu t0, t1, 2f
     sb zero, 0(t0)
     addi t0, t0, 1
     j 1b
 2:
     call machine_init
\end{verbatim}
\end{codesnippet}

Detalles técnicos:

\begin{itemize}
  \item \textbf{\texttt{gp} y \texttt{.option norelax}.} La psABI reserva \texttt{gp} para accesos \emph{small-data}
  gp-relativos~\cite{riscv-psabi}. Si no se desactiva la \emph{linker relaxation} alrededor de \texttt{la gp,
  \_\_global\_pointer\$}, el linker ve un \texttt{auipc}/\texttt{addi} y lo reescribe como un \texttt{addi gp, gp, 0}
  gp-relativo --- básicamente, asume que \texttt{gp} ya vale lo que tiene que valer. Resultado: \texttt{gp} queda en
  basura y cualquier acceso small-data posterior lee cualquier cosa. \texttt{.option norelax} apaga esa optimización en
  ese bloque puntual.

  \item \textbf{Stack del kernel.} Es el equivalente a ``setear el stack en \texttt{0x27000}'' del TP x86. El símbolo
  \texttt{kernel\_stack\_top} está definido al final de un bloque en \texttt{.bss} de 8~KiB, resultado del linker (ver
  \texttt{riscv/linker.ld}). No hace falta una dirección mágica porque el linker coloca todo a partir de
  \texttt{0x80000000} y la ubicación concreta de el stack la decide el linker al resolver símbolos.
\end{itemize}

\subsection{\texttt{riscv/machine.S}: mini firmware en M-mode}
\label{sec:machine}

\texttt{machine\_init} cumple el rol que en un sistema típico cumple OpenSBI (la firmware SBI de referencia en
RISC-V~\cite{riscv-sbi}). Como se corre con \texttt{-bios none}, no hay SBI, y el kernel tiene que proveer él mismo la
configuración mínima de M-mode. El esqueleto sigue la receta \emph{Bare Bones} para RISC-V documentada en
OSDev~\cite{osdev-riscv}. Los pasos:

\begin{enumerate}
  \item \textbf{Stack y trap vector de M-mode.} \texttt{mscratch} guarda el tope de el stack de M-mode y \texttt{mtvec}
  apunta a \texttt{machine\_trap\_entry}.
  \item \textbf{Delegación.} \texttt{mideleg} delega a S-mode los interrupts de supervisor (SSIP, STIP, SEIP) y
  \texttt{medeleg} delega las excepciones estándar 0..15. Así, las excepciones desde U-mode saltan directamente a
  \texttt{stvec} del kernel y no a \texttt{mtvec}.
  \item \textbf{PMP.} Una sola región TOR que abarca todo el espacio (\texttt{pmpaddr0 = -1}, \texttt{pmpcfg0 =
  RWX+TOR}) le da a S-mode y U-mode acceso pleno según PMP. El control fino después lo dan las PTEs Sv39.
  \item \textbf{Timer.} Se habilita MTIE en \texttt{mie} y se programa el primer \texttt{mtimecmp} = \texttt{time} +
  \texttt{g\_timer\_delta\_ticks}.
  \item \textbf{Salto a S-mode.} Se setea \texttt{mstatus.MPP}$=$1 (S), \texttt{mstatus.MPIE}$=$1,
  \texttt{mepc}$=$\texttt{kernel\_main} y se ejecuta \texttt{mret}. Esto es el análogo RISC-V del \texttt{jmp
  0xA0:modo\_protegido} del TP x86.
\end{enumerate}

\begin{codesnippet}
\begin{verbatim}
 # Delegar interrupts a S-mode: SSIP, STIP, SEIP
 li t0, ((1 << 1) | (1 << 5) | (1 << 9))
 csrw mideleg, t0

 # Delegar excepciones 0..15 a S-mode
 li t0, 0xFFFF
 csrw medeleg, t0

 # PMP: R|W|X + TOR cubriendo 0..max
 li t0, -1
 csrw pmpaddr0, t0
 li t0, 0x0F
 csrw pmpcfg0, t0
\end{verbatim}
\end{codesnippet}

\subsection{``Segmento de pantalla'' y limpieza inicial}

Como no hay segmentación, la pantalla se resuelve por \textbf{mapeo de MMIO}: el UART 16550 de QEMU \texttt{virt} en
\texttt{0x1000\_0000}. En cada page-table de tarea (y en la de idle) se mapea ese rango con una megapage (2~MiB) a nivel
L1, con flags \texttt{V|R|W|A|D}, sin \texttt{U}, para que el kernel pueda escribirlo. El equivalente al ``segmento de
pantalla sólo accesible por el kernel'' del TP x86 es entonces el bit \texttt{U}$=$0 de esa entrada.

La limpieza y pintado del tablero equivalen, a nivel código, a \texttt{screen\_draw} / \texttt{screen\_drawBox} en
\texttt{riscv/screen.cpp}. Internamente la pantalla se emula como un buffer \texttt{g\_video[50][80]} de celdas
(caracter + atributo estilo VGA), que se vuelca al UART mediante secuencias ANSI (\texttt{\textbackslash x1b[?1049h}
para el alternate screen, \texttt{\textbackslash x1b[y;xH} para posicionar el cursor, \texttt{\textbackslash x1b[fg;bgm}
para atributos, etc.). Al imprimir, un atributo estilo \texttt{C\_BG\_RED | C\_FG\_WHITE} se traduce a los códigos ANSI
\texttt{41} y \texttt{97}.

\newpage
\section{Ejercicio 2 --- Manejo de excepciones: IDT $\to$ \texttt{stvec}}
\subsection{Qué hacía el TP x86}

Se armaba una IDT de 256 entradas. Para cada excepción (0..18) se daba de alta una \texttt{IDT\_ENTRY} con el selector
de código de kernel, DPL=0, Present=1, Type=1110 (Interrupt Gate). Cada ISR (\_is0, \_is1, \ldots) se generaba con una
macro en \texttt{isr.asm} e imprimía un mensaje. Se cargaba la IDT con \texttt{lidt}.

\subsection{Equivalente en RISC-V}

RISC-V no tiene una tabla indexada de gates: hay un único CSR, \texttt{stvec}, que apunta a la dirección a la que salta
el procesador siempre que una excepción o interrupción tiene que atenderse en S-mode. El motivo lo decide el handler en
software leyendo \texttt{scause}. Con lo cual:

\begin{itemize}
  \item La IDT se reemplaza por una \textbf{rutina única en ASM} (\texttt{trap\_entry}, en \texttt{riscv/trap.S}) que
  guarda todo el contexto y llama a un despachador en C++ (\texttt{trap\_handle}, en \texttt{riscv/kernel.cpp}).
  \item ``Selector + offset + DPL'' ya no tiene sentido: \texttt{stvec} es una dirección, y el privilegio al que vuelve
  el handler sale de \texttt{sstatus.SPP}.
  \item El ``cargar IDT con \texttt{lidt}'' queda en:
\begin{codesnippet}
\begin{verbatim}
 asm volatile("csrw stvec, %0" : : "r"(trap_entry));
\end{verbatim}
\end{codesnippet}
  en \texttt{kernel\_main}, con \texttt{stvec} en modo ``Direct'' (MODE=0). El modo ``Vectored'' (MODE=1) permitiría un
  salto por causa de interrupción; no se usa aquí.
\end{itemize}

\subsection{Guardado de contexto: \texttt{TrapFrame} y \texttt{trap.S}}

El equivalente de ``\texttt{pushad} + selectores'' es la estructura \texttt{TrapFrame} (\texttt{riscv/trap.h}):

\begin{codesnippet}
\begin{verbatim}
 struct TrapFrame {
   uint64_t ra, sp, gp, tp;
   uint64_t t0, t1, t2;
   uint64_t s0, s1;
   uint64_t a0, a1, a2, a3, a4, a5, a6, a7;
   uint64_t s2, ..., s11;
   uint64_t t3, t4, t5, t6;
   uint64_t sepc, sstatus, scause, stval;
 };
\end{verbatim}
\end{codesnippet}

\texttt{trap\_entry} (en \texttt{riscv/trap.S}) hace, a grandes rasgos:

\begin{enumerate}
  \item Cambia a una stack de trap dedicada (\texttt{trap\_stack\_top}). Esto es necesario porque, al entrar desde
  U-mode, \texttt{sp} apunta al stack del usuario, y usarlo implicaría ejecutar código de kernel en una página de
  usuario.
  \item Guarda todos los GPRs (excepto \texttt{x0}) en el \texttt{TrapFrame} apuntado por
  \texttt{g\_current\_trapframe}.
  \item \textbf{Recarga \texttt{gp} al del kernel} con \texttt{.option norelax}. Esta parte es crítica: los binarios de
  usuario se linkean en \texttt{0x08000000} y tienen su propio \texttt{\_\_global\_pointer\$} $\sim$\texttt{0x08000800},
  por lo que los accesos gp-relativos del kernel (pensados con el gp del kernel, $\sim$\texttt{0x8000....}) faltarían si
  el trap handler mantuviera el gp del usuario.
  \item Guarda \texttt{sepc}, \texttt{sstatus}, \texttt{scause}, \texttt{stval}.
  \item Llama a \texttt{trap\_handle(tf)}, que devuelve un nuevo \texttt{TrapFrame*} (la ``próxima tarea'' a correr).
  \item Restaura el contexto del \texttt{TrapFrame} retornado (incluyendo \texttt{sepc}/\texttt{sstatus}) y ejecuta
  \texttt{sret}.
\end{enumerate}

La rutina \texttt{enter\_user}, también en \texttt{trap.S}, se utiliza para el salto inicial desde \texttt{kernel\_main}
a la primera tarea U-mode: setea \texttt{g\_current\_trapframe}, carga directamente todos los registros desde el
\texttt{TrapFrame} y ejecuta \texttt{sret}. Es el análogo RISC-V del \texttt{jmp 0xB8:0} del TP x86 para saltar a la
tarea idle.

\subsection{Probar una excepción}

El TP x86 probaba la IDT forzando una división por cero. En este port, el binario de usuario \texttt{riscv/user/miner.c}
provoca deliberadamente una excepción cuando el minero tiene id 5:

\begin{codesnippet}
\begin{verbatim}
 if (id == 5) {
     syscall_move(Forward);  /* ... */
     *(volatile uint32_t*)0x0 = 0xDEADBEEF;
     for (;;) {}
 }
\end{verbatim}
\end{codesnippet}

Escribir en la dirección virtual \texttt{0x0} desde U-mode genera un \emph{Store/AMO page fault} (\texttt{scause}$=$15)
porque esa dirección no está mapeada. \texttt{trap\_handle} mata a la tarea mediante \texttt{sched\_kill()} y cede el
quantum a idle (ver ejercicio~7).

\newpage
\section{Ejercicio 3 --- Reloj, teclado e interrupción 47}
\subsection{Reloj: PIT + PIC $\to$ MTIP + SSIP}
\label{sec:reloj}

\subsubsection*{Qué hacía el TP x86}

La \texttt{\_isr32} llamaba a \texttt{nextClock} (animación del cursor) y \texttt{pic\_finish1} (EOI al PIC maestro), y
ejecutaba \texttt{iret}.

\subsubsection*{Equivalente en RISC-V}

En RISC-V el timer está en CLINT/ACLINT~\cite{riscv-aclint}: dos MMIOs, \texttt{mtime} (contador libre) y
\texttt{mtimecmp} (comparador por HART). Cuando \texttt{mtime} $\geq$ \texttt{mtimecmp}, se levanta \textbf{MTIP}
(machine timer interrupt pending). El problema es que MTIP \textbf{solo} se puede atender en M-mode (no se puede delegar
a S-mode directamente~\cite{riscv-priv}). Un sistema con OpenSBI resolvería esto usando
\texttt{sbi\_set\_timer}~\cite{riscv-sbi}; con \texttt{-bios none} no hay SBI y debemos resolverlo nosotros.

El patrón estándar (el mismo que usa xv6 en su build sin SBI~\cite{xv6-riscv} y que documenta el wiki OSDev para
RISC-V~\cite{osdev-riscv}) es:

\begin{enumerate}
  \item Atender el MTIP en M-mode (\texttt{machine\_trap\_entry}, \texttt{mcause}$=$7).
  \item Reprogramar el próximo \texttt{mtimecmp}.
  \item \textbf{Levantar un supervisor software interrupt (SSIP)} escribiendo \texttt{sip.SSIP}$=$1 desde M-mode.
  \item Ejecutar \texttt{mret} y dejar que la CPU dispare el handler S-mode (vía \texttt{stvec}) con \texttt{scause}$=$1
  (SSIP).
\end{enumerate}

En \texttt{machine.S}:

\begin{codesnippet}
\begin{verbatim}
 # Reprogramar mtimecmp
 csrr t0, time
 la t1, g_timer_delta_ticks
 ld t1, 0(t1)
 add t0, t0, t1
 csrr t1, mhartid
 slli t1, t1, 3
 li t2, MTIMECMP_BASE        # 0x0200_4000 + 8*hartid
 add t2, t2, t1
 sd t0, 0(t2)

 # Levantar SSIP
 li t0, 2
 csrs sip, t0
\end{verbatim}
\end{codesnippet}

Y del lado de S-mode, \texttt{trap\_handle} reconoce SSIP como el tick:

\begin{codesnippet}
\begin{verbatim}
 if (is_interrupt) {
     if (cause_code == 1) {  // SSIP
         const uint64_t ssip_mask = (1ull << 1);
         asm volatile("csrc sip, %0" : : "r"(ssip_mask));
         /* cooldowns + poll UART + spinner + scheduler */
         return schedule_next();
     }
     return tf;
 }
\end{verbatim}
\end{codesnippet}

El \texttt{csrc sip, SSIP} es el análogo de \texttt{pic\_finish1}: limpia el bit pendiente para que el interrupt no se
vuelva a disparar. La función \texttt{nextClock} del TP original (animación del spinner en la esquina) se preserva con
\texttt{update\_spinner()} en \texttt{kernel.cpp}.

\subsection{Teclado: IRQ1 + scancodes $\to$ UART polling}

\subsubsection*{Qué hacía el TP x86}

La \texttt{\_isr33} leía el scancode de \texttt{in al, 0x60}, pusheaba \texttt{eax} y llamaba a
\texttt{idt\_int\_teclado} (que traducía el scancode a la acción del juego), luego \texttt{pic\_finish1} y
\texttt{iret}.

\subsubsection*{Equivalente en RISC-V}

QEMU \texttt{virt} no tiene teclado PS/2 estándar; lo que hay es el UART 16550 sobre \texttt{-serial stdio}. Podría
tomarse la interrupción del UART vía PLIC, pero para no sumar un mecanismo más al informe terminé haciéndolo por
\textbf{polling en cada tick}, que a la frecuencia del reloj ($\sim$20~Hz) se percibe igual. En lugar de un ISR por
tecla, \texttt{trap\_handle} llama a \texttt{poll\_uart\_input()} al atender el tick, y ese drena el buffer con
\texttt{uart::getc\_nonblock()}:

\begin{codesnippet}
\begin{verbatim}
 int uart::getc_nonblock() {
     if ((mmio_read8(UART_LSR) & LSR_DATA_READY) == 0)
         return -1;
     return mmio_read8(UART_RHR);
 }
\end{verbatim}
\end{codesnippet}

Al llegar un byte se lo pasa a \texttt{handle\_key}, que mapea teclas (no scancodes) a acciones de juego:

\begin{itemize}
  \item Jugador ROJO: \texttt{W} (arriba), \texttt{S} (abajo), \texttt{Q}/\texttt{E} (spawnear minero).
  \item Jugador AZUL: \texttt{I} (arriba), \texttt{K} (abajo), \texttt{O}/\texttt{P} (spawnear minero).
\end{itemize}

El byte recibido se imprime en la esquina superior derecha en hex (equivalente al ``mostrar la tecla 0..9 en pantalla''
del TP x86), útil para validar input.

\subsection{\texttt{int 47} (x86) $\to$ \texttt{ecall} (RISC-V)}

En x86 se declaraba una Interrupt Gate con DPL=3 en la entrada 47 de la IDT, y la \texttt{\_isr47} inicialmente seteaba
\texttt{eax}$=$\texttt{0x42} y hacía \texttt{iret}. En RISC-V el análogo de esa \emph{software interrupt} es
\texttt{ecall}: desde U-mode genera un trap con \texttt{scause}$=$8 (``Environment call from U-mode''). No hay un
``número'' configurable como en \texttt{int N}; lo que sí se mantiene del TP3 son los IDs de servicio, que en este port
van por \texttt{a7} siguiendo la convención estándar de syscalls de RISC-V:

\begin{center}
\begin{tabular}{l|l}
\textbf{Servicio} & \textbf{ABI} \\
\hline
\texttt{move(d)} & \texttt{a7}=177788, \texttt{a0}=\texttt{d} \\
\texttt{take()}  & \texttt{a7}=310311, retorna en \texttt{a0} \\
\texttt{getId()} & \texttt{a7}=760279, retorna en \texttt{a0} \\
\end{tabular}
\end{center}

No se usa un ``eax transitorio con \texttt{0x42}'' porque el handler puede escribir directamente sobre \texttt{tf->a0} y
el \texttt{sret} restaura \texttt{a0} como valor de retorno. El stub de usuario vive en \texttt{riscv/syscall.h}:

\begin{codesnippet}
\begin{verbatim}
 static inline int32_t syscall_take() {
     register uint64_t a7 asm("a7") = 310311;
     register uint64_t a0 asm("a0");
     asm volatile("ecall" : "=r"(a0) : "r"(a7) : "memory");
     return (int32_t)a0;
 }
\end{verbatim}
\end{codesnippet}

El despacho en \texttt{trap\_handle} se describe en el ejercicio~7.

\newpage
\section{Ejercicio 4 --- Paginación del kernel}
\subsection{Qué hacía el TP x86}

\texttt{mmu\_initKernelDir} armaba un Page Directory en \texttt{0x27000} y una PT en \texttt{0x28000}, con identity
mapping de los primeros 4~MiB. Luego \texttt{kernel.asm} cargaba CR3 con \texttt{KERNEL\_PAGE\_DIR} y seteaba
\texttt{CR0.PG}$=$1.

\subsection{Equivalente en RISC-V: Sv39}
\label{sec:sv39}

Sv39 es el esquema de paginación de 3 niveles más común en rv64:

\begin{itemize}
  \item Direcciones virtuales de 39 bits (los superiores deben ser sign-extended).
  \item Tres niveles de page table con 512 entradas de 8 bytes cada una (4~KiB por tabla). VPN2 (bits 38:30), VPN1 (bits
  29:21), VPN0 (bits 20:12).
  \item Páginas de 4~KiB, 2~MiB (\emph{megapage}, leaf en L1) o 1~GiB (\emph{gigapage}, leaf en L2).
  \item PTE de 64 bits con los bits \texttt{V} (1), \texttt{R} (2), \texttt{W} (4), \texttt{X} (8), \texttt{U} (16),
  \texttt{G} (32), \texttt{A} (64), \texttt{D} (128), seguidos de la PPN.
\end{itemize}

El equivalente conjunto de ``\texttt{mov cr3, eax} + \texttt{or CR0, PG}'' es escribir el CSR \texttt{satp} con formato:
\[
\texttt{satp} = (\underbrace{\texttt{MODE}}_{\text{SV39}=8} \ll 60) \;|\; (\texttt{ASID} \ll 44) \;|\;
\texttt{PPN}(\text{root})
\]
seguido de un \texttt{sfence.vma} para invalidar caches de traducción. En \texttt{mmu.cpp}:

\begin{codesnippet}
\begin{verbatim}
 static inline uint64_t make_satp(uint64_t root_pa, uint16_t asid) {
     const uint64_t ppn = root_pa >> 12;
     return (SATP_MODE_SV39 << 60)
          | (static_cast<uint64_t>(asid) << 44)
          | (ppn & 0x0FFF'FFFF'FFFFull);
 }
\end{verbatim}
\end{codesnippet}

\subsection{Identity mapping del kernel}

Dado que en RISC-V rv64 el espacio virtual es holgado, en lugar de armar una PT nivel-0 que mapee los primeros 4~MiB
como en el TP x86, se usa una única \textbf{leaf de 1~GiB en L2} para cubrir todo el rango de RAM relevante
(\texttt{0x8000\_0000}..\texttt{0xBFFF\_FFFF}), con flags \texttt{V|R|W|X|A|D} (supervisor only, sin \texttt{U}):

\begin{codesnippet}
\begin{verbatim}
 static inline void map_kernel_gigapage(uint64_t* root) {
     const uint64_t va = 0x8000'0000ull;
     const uint64_t pa = 0x8000'0000ull;
     const uint64_t flags = PTE_V|PTE_R|PTE_W|PTE_X|PTE_A|PTE_D;
     root[vpn2(va)] = pte_make_leaf(pa, flags);
 }
\end{verbatim}
\end{codesnippet}

Esto reemplaza la tabla \texttt{KERNEL\_PAGE\_TABLE\_0} del TP x86. Los bits \texttt{A} y \texttt{D} son requeridos por
el hardware en el momento del primer acceso; setearlos \emph{a priori} simplifica la implementación y está
explícitamente permitido por la especificación~\cite{riscv-priv}.

\subsection{UART como ``MMIO del video''}

El MMIO del UART 16550 se mapea como una \textbf{megapage} (2~MiB) en L1 con flags \texttt{V|R|W|A|D} (sin \texttt{X} ni
\texttt{U}). Este bloque cumple el rol que en x86 cumplía el ``segmento de pantalla''. El mapeo está presente tanto en
la page-table del kernel como en la de cada tarea (aunque las tareas no tienen \texttt{U} en esa entrada, por lo que no
pueden tocar el UART desde modo usuario).

\subsection{Activación de paginación y PMP}

Antes del \texttt{satp}, PMP tiene que autorizar el acceso a memoria en S-mode. Eso ya lo hizo \texttt{machine\_init} en
M-mode (sección~\ref{sec:machine}). Por lo tanto, la ``activación de paginación'' del TP3 se reduce a una sola línea en
\texttt{mmu\_init}:

\begin{codesnippet}
\begin{verbatim}
 void mmu_init() {
     if (!g_idle.initialized) { /* ... construye g_idle ... */ }
     mmu_switch_to_idle();   // csrw satp + sfence.vma
 }
\end{verbatim}
\end{codesnippet}

A partir de esa llamada, todas las direcciones en S-mode son traducidas por Sv39.

\newpage
\section{Ejercicio 5 --- MMU de tareas}
\subsection{Qué hacía el TP x86}

Este ejercicio era el más denso del TP3. Incluía:

\begin{itemize}
  \item \texttt{mmu\_init}: contadores \texttt{next\_free\_kernel\_page} (desde \texttt{0x100000}) y
  \texttt{next\_free\_task\_page} (desde \texttt{0x400000}).
  \item \texttt{mmu\_initTaskDir}: construir el PD/PT de la tarea, identity-mapear kernel + área libre kernel, copiar el
  código original de la tarea (desde \texttt{MMU\_BASE\_ADDR\_TAREAS}) a su celda física en el tablero, mapear
  \texttt{TASK\_CODE} (\texttt{0x08000000}) a esas páginas con permisos U, y mapear \texttt{TASK\_COMPARTIDA} a
  \texttt{g\_shared[jugador]}.
  \item \texttt{mmu\_mapPage} / \texttt{mmu\_unmapPage}: helpers para mapear/desmapear una página virtual, invalidando
  el TLB con \texttt{tlbflush}.
\end{itemize}

\subsection{Equivalente en RISC-V}

La portación conserva el esquema ``el código del minero vive físicamente en su celda del tablero''. Las funciones
equivalentes viven en \texttt{riscv/mmu.cpp}:

\begin{itemize}
  \item \texttt{mmu\_task\_init(slot, player, x, y)}: equivalente de \texttt{mmu\_initTaskDir}.
  \item \texttt{mmu\_task\_move(slot, old\_x, old\_y, new\_x, new\_y)}: equivalente conjunto de ``copiar 8~KiB +
  remapear \texttt{TASK\_CODE}''.
  \item \texttt{mmu\_switch\_to\_task(slot)} / \texttt{mmu\_switch\_to\_idle()}: análogo a ``cargar CR3'' (escribir
  \texttt{satp}).
\end{itemize}

No se exponen \texttt{mmu\_mapPage}/\texttt{mmu\_unmapPage} genéricas, porque en este port el universo de mapeos es más
restringido y conocido (kernel gigapage, UART megapage, dos páginas de código y dos páginas compartidas por tarea). El
rol de esas funciones lo cumplen \texttt{map\_user\_code\_pages} y \texttt{map\_user\_shared}, que actualizan
directamente \texttt{t.l0\_user[]}.

\subsection{Allocator del kernel}

El análogo del ``contador de páginas libres'' del TP x86 es un \emph{bump allocator} a partir del final del
\texttt{.bss}:

\begin{codesnippet}
\begin{verbatim}
 static uint64_t g_alloc_page = 0;

 static void* alloc_page_table() {
     if (g_alloc_page == 0) {
         g_alloc_page = align_up_u64(
             reinterpret_cast<uint64_t>(&__bss_end), PAGE_SIZE);
     }
     void* page = reinterpret_cast<void*>(g_alloc_page);
     g_alloc_page += PAGE_SIZE;
     memset(page, 0, PAGE_SIZE);
     return page;
 }
\end{verbatim}
\end{codesnippet}

\subsection{``Tablero como memoria''}

El ``área mapa'' del TP3 (\texttt{0x400000}..\texttt{0x1C5FFFF}, 78$\times$40 celdas de 8~KiB) se reemplaza por un
\textbf{arreglo estático alineado a página}:

\begin{codesnippet}
\begin{verbatim}
 alignas(PAGE_SIZE) static uint8_t
     g_board[TABLERO_N_COLS * TABLERO_N_FILS * kCellBytes];

 static inline uint64_t board_cell_pa(uint32_t x, uint32_t y) {
     const uint64_t ix = y * TABLERO_N_COLS + x;
     return reinterpret_cast<uint64_t>(&g_board[ix * kCellBytes]);
 }
\end{verbatim}
\end{codesnippet}

Esto evita tener que fijar direcciones físicas ``mágicas'' y deja que el linker resuelva el layout. Igualmente se cumple
la propiedad del TP3: cada celda del tablero es un bloque contiguo de 8~KiB y \texttt{board\_cell\_pa(x,y)} es su
dirección física.

Las áreas compartidas entre mineros del mismo jugador se resuelven también con arreglos estáticos:

\begin{codesnippet}
\begin{verbatim}
 alignas(PAGE_SIZE) static uint8_t g_shared[2][2 * PAGE_SIZE];
\end{verbatim}
\end{codesnippet}

\subsection{Construcción de la page-table de tarea}

\texttt{mmu\_task\_init} construye la page-table completa de una tarea:

\begin{codesnippet}
\begin{verbatim}
 static inline void task_build_pagetable(TaskMmu& t,
         uint16_t asid, uint32_t player_id, uint64_t cell_pa) {
     t.root = reinterpret_cast<uint64_t*>(alloc_page_table());
     map_kernel_gigapage(t.root);
     ensure_l1_0(t);
     map_uart(t);
     map_user_code_pages(t, cell_pa);       // TASK_CODE (2 pags)
     map_user_shared(t, player_id);         // TASK_COMPARTIDA (2 pags)
     t.asid = asid;
     t.satp = make_satp(
         reinterpret_cast<uint64_t>(t.root), asid);
     t.initialized = 1;
 }
\end{verbatim}
\end{codesnippet}

El mapeo de \texttt{TASK\_CODE} es el punto crítico de la mecánica del juego: a partir de la VA \texttt{0x08000000} se
mapean dos páginas físicas contiguas a partir de \texttt{cell\_pa}, con flags \texttt{V|R|W|X|U|A|D}. Es decir, \emph{el
minero puede leer, escribir y ejecutar en el mismo bloque de memoria} --- exactamente la condición que hace posible que
el juego copie el código de la tarea a otra celda y la tarea siga ejecutando desde \texttt{0x08000000} con código
distinto.

\subsection{\texttt{mmu\_task\_move}: movimiento del minero}

Es el análogo a ``copiar 8~KiB + \texttt{mmu\_mapPage(TASK\_CODE, ...)} + \texttt{tlbflush}'' del TP x86:

\begin{codesnippet}
\begin{verbatim}
 void mmu_task_move(uint32_t slot,
                    uint32_t old_x, uint32_t old_y,
                    uint32_t new_x, uint32_t new_y) {
     TaskMmu& t = g_tasks[slot];
     if (!t.initialized) return;

     const uint64_t src_pa = board_cell_pa(old_x, old_y);
     const uint64_t dst_pa = board_cell_pa(new_x, new_y);

     memcpy((void*)dst_pa, (const void*)src_pa, kCellBytes);
     fence_i();                              // ifetch coherence

     map_user_code_pages(t, dst_pa);
     sfence_vma_addr(kUserBase);
     sfence_vma_addr(kUserBase + PAGE_SIZE);
 }
\end{verbatim}
\end{codesnippet}

Detalles técnicos:

\begin{itemize}
  \item \textbf{\texttt{fence.i}.} Después de escribir instrucciones en memoria, RISC-V no garantiza que el pipeline de
  instrucciones vea los bytes nuevos hasta una \texttt{fence.i} explícita~\cite{riscv-unpriv}. En x86 la coherencia
  I\$/D\$ es automática; en RISC-V no, y sin esta instrucción el minero podría ejecutar el código viejo después del
  movimiento (me pasó debuggeando: la tarea arrancaba, se ``movía'', y seguía ejecutando como si estuviera en la celda
  original).
  \item \textbf{\texttt{sfence.vma rs1=va, rs2=x0}.} Invalida sólo la entrada del TLB para esa VA, el equivalente del
  \texttt{invlpg} de x86. Es más barato que un flush global y alcanza porque el único cambio fue \texttt{TASK\_CODE}.
\end{itemize}

\subsection{ASIDs}

Sv39 permite incluir un \texttt{ASID} en \texttt{satp} para que el TLB no haga flush global al cambiar de address space.
Este port asigna:

\begin{itemize}
  \item \texttt{asid}$=$0 para idle.
  \item \texttt{asid}$=$slot+1 para mineros (1..20).
\end{itemize}

El \texttt{sfence.vma x0, x0} sólo se usa al inicializar una tarea; en runtime los cambios de address space se hacen con
\texttt{mmu\_switch()} que escribe \texttt{satp} (con su ASID) y hace un \texttt{sfence.vma} global (simplicidad $>$
eficiencia).

\newpage
\section{Ejercicio 6 --- Contexto de tarea: TSS $\to$ TrapFrame}
\subsection{Qué hacía el TP x86}

El TP3 usaba \textbf{task switching por hardware}: un descriptor de TSS en la GDT por tarea
(\texttt{tss\_initGdtEntry}), una TSS por tarea (\texttt{fill\_tss\_idle}, \texttt{fill\_tss}) con \texttt{eip},
\texttt{cr3}, \texttt{esp}, \texttt{eflags}, selectores de segmento, etc. El cambio de tarea se hacía con un \texttt{jmp
<sel>:<off>} al selector de la TSS, dejando que el procesador guardara/cargara todo el contexto automáticamente.

\subsection{Equivalente en RISC-V}

RISC-V no tiene task switching por hardware, todo es por software. Se simplifica bastante: no hay GDT, ni TSS, ni
descriptores; queda una estructura por tarea y nada más:

\begin{codesnippet}
\begin{verbatim}
 struct Task {
   TrapFrame tf;
   uint32_t initialized;
 };
 static Task g_tasks[20];
 static TrapFrame g_idle_tf;
\end{verbatim}
\end{codesnippet}

Cada \texttt{TrapFrame} cumple el rol que en el TP x86 cumplía la TSS: contiene el \emph{snapshot} completo del estado
de esa tarea. El campo \texttt{initialized} replica la noción de ``descriptor presente''.

\subsection{Inicialización de \texttt{TrapFrame} (equivalente de \texttt{fill\_tss})}

Toda la ``configuración inicial de la TSS'' del TP x86 (eip, esp, cs/ds/ss, eflags) se colapsa en tres campos del
\texttt{TrapFrame}:

\begin{codesnippet}
\begin{verbatim}
 static void init_user_trapframe(TrapFrame& tf, uint64_t entry_pc) {
     tf = {};
     tf.sepc = entry_pc;              // PC inicial
     tf.sp   = TASK_STACK;            // stack de usuario
     // SPIE=1 (SIE <- SPIE tras sret), SPP=0 (retorno a U-mode).
     tf.sstatus = (1ull << 5);
 }
\end{verbatim}
\end{codesnippet}

\begin{itemize}
  \item \texttt{sepc}$=$\texttt{TASK\_CODE} (\texttt{0x08000000}) $\Leftrightarrow$ \texttt{tss.eip}$=$dirección inicial
  del código.
  \item \texttt{sp}$=$\texttt{TASK\_STACK}$=$\texttt{TASK\_CODE}+2$\cdot$4~KiB $\Leftrightarrow$ \texttt{tss.esp}.
  \item \texttt{sstatus.SPIE}$=$1 asegura que \texttt{SIE}$=$1 luego del \texttt{sret}, es decir, la tarea corre con
  interrupciones habilitadas (equivalente a \texttt{eflags.IF}$=$1 del TP x86, donde se usaba \texttt{0x202}).
  \item \texttt{sstatus.SPP}$=$0 indica que el \texttt{sret} devuelve el control a U-mode (equivalente a ``DPL=3''). No
  hay ``CS'' ni ``DS''; el modo es el único eje de privilegio.
\end{itemize}

Nótese que en RISC-V no existen ``registros de segmento'': eso elimina por completo los campos
\texttt{cs/ds/ss/fs/gs/es} de la TSS. Tampoco hace falta una ``stack de nivel 0'' por tarea (\texttt{ss0:esp0}): al
entrar al kernel vía trap, se cambia a el stack de trap única (\texttt{trap\_stack\_top}), que alcanza porque SIE se
limpia al entrar al trap (no hay traps anidados al mismo nivel de privilegio).

\subsection{``Salto a la tarea Idle''}

El TP x86 arrancaba con \texttt{jmp 0xB8:0} después de cargar \texttt{tr} con el selector de la tarea inicial. En este
port, \texttt{kernel\_main} hace:

\begin{codesnippet}
\begin{verbatim}
 init_idle_once();

 asm volatile("csrw stvec, %0" : : "r"((uint64_t)trap_entry));
 const uint64_t sie_mask = (1ull << 1) | (1ull << 5); // SSIE + STIE
 asm volatile("csrs sie, %0" : : "r"(sie_mask));

 TrapFrame* first = schedule_next();
 enter_user(first);
\end{verbatim}
\end{codesnippet}

\texttt{enter\_user} (en \texttt{trap.S}) restaura todos los GPRs del \texttt{TrapFrame} y ejecuta \texttt{sret}. Con
\texttt{sstatus.SPP}$=$0 y \texttt{sepc}$=$\texttt{TASK\_CODE}, esto transiciona a U-mode saltando al primer byte del
binario de idle.

\subsection{La tarea Idle}

En el TP x86, idle era un binario del kernel en \texttt{0x00014000} que compartía \texttt{cr3} con el kernel. En este
port idle es un binario de usuario \emph{separado} (\texttt{riscv/user/idle.c}, ver sección~\ref{sec:user}), con su
propia address space (\texttt{asid}$=$0) y su propio buffer de 8~KiB como ``celda'' (no está en el tablero). Su código
completo es:

\begin{codesnippet}
\begin{verbatim}
 void user_main() {
     for (;;) __asm__ volatile("nop");
 }
\end{verbatim}
\end{codesnippet}

El porqué de hacerla U-mode y no S-mode: mantiene el mismo ABI/path de entrada/salida que los mineros, simplifica el
\texttt{TrapFrame} y evita tener que distinguir ``contexto supervisor'' vs ``contexto usuario'' en el scheduler.

\newpage
\section{Ejercicio 7 --- Scheduler y servicios}
\subsection{Estructuras del scheduler}

Las estructuras en \texttt{riscv/state.h} son un calco directo del TP3:

\begin{codesnippet}
\begin{verbatim}
 sched_miner_state_t sched_miners_state[20];
 uint32_t sched_player_row[2];
 uint32_t sched_miners_left[2];
 uint8_t  sched_last_miner[2];
 uint8_t  sched_this_miner;
\end{verbatim}
\end{codesnippet}

La inicialización (\texttt{sched\_init}) hace exactamente lo mismo que la versión x86: marca los 20 slots como muertos,
reinicia contadores y filas. La única diferencia relevante es que, en este port, no hay una ``GDT de TSS'' que resetear:
los slots son simples índices.

\subsection{\texttt{sched\_nextTask}}

El algoritmo del TP3 (alternar entre jugadores y, si el opuesto no tiene mineros vivos, usar uno del jugador actual; si
ninguno hay, usar idle) se preserva:

\begin{codesnippet}
\begin{verbatim}
 int32_t sched_nextTask() {
     uint32_t actual   = (sched_this_miner < 10) ? PLAYER_A : PLAYER_B;
     uint32_t opuesto  = (actual == PLAYER_A) ? PLAYER_B : PLAYER_A;
     uint16_t ix = 0;

     if (aux_get_next_alive_miner(opuesto, &ix)) {
         sched_this_miner = ix;
         sched_last_miner[opuesto] = ix % 10;
         return ix;
     }
     if (aux_get_next_alive_miner(actual, &ix)) {
         sched_this_miner = ix;
         sched_last_miner[actual] = ix % 10;
         return ix;
     }
     return -1;  // idle
 }
\end{verbatim}
\end{codesnippet}

La diferencia clave con el informe x86 es el \emph{tipo de retorno}: ya no es un selector de GDT, sino el índice 0..19
del minero (o $-1$ para idle). Quien consume ese índice es \texttt{schedule\_next} en \texttt{kernel.cpp}, que convierte
el slot a \texttt{TrapFrame*} y hace \texttt{mmu\_switch\_to\_task(slot)}:

\begin{codesnippet}
\begin{verbatim}
 static TrapFrame* schedule_next() {
     int32_t next = sched_nextTask();
     if (next < 0) return yield_to_idle();
     uint32_t slot = (uint32_t)next;
     init_task_if_needed(slot);
     mmu_switch_to_task(slot);
     g_current_slot = next;
     sched_this_miner = (uint8_t)next;
     return &g_tasks[slot].tf;
 }
\end{verbatim}
\end{codesnippet}

\subsection{Despacho de \texttt{ecall}}

La ``rutina de int 47'' del TP3 se vuelve un \emph{branch} dentro de \texttt{trap\_handle} cuando \texttt{scause}$=$8:

\begin{codesnippet}
\begin{verbatim}
 if (cause_code == 8) {            // ecall desde U-mode
     tf->sepc += 4;                // saltar la propia ecall
     poll_uart_input();
     const uint64_t sys_id = tf->a7;

     if (sys_id == 177788) {       // move(d=a0)
         game_move((e_direction_t)tf->a0);
         return yield_to_idle();   // equiv. del "jmp idle" del TP
     }
     if (sys_id == 310311) {       // take() -> a0
         tf->a0 = (uint64_t)game_take();
         return yield_to_idle();
     }
     if (sys_id == 760279) {       // getId() -> a0
         tf->a0 = (uint64_t)game_get_id();
         return tf;                // sin yield
     }
     if (g_current_slot >= 0) sched_kill();
     return yield_to_idle();
 }
\end{verbatim}
\end{codesnippet}

Detalles:

\begin{itemize}
  \item \texttt{sepc += 4}: a diferencia de \texttt{int N} en x86 (que pushea la dirección \emph{siguiente}), en RISC-V
  \texttt{sepc} apunta a la \texttt{ecall} misma. Hay que avanzarlo manualmente al retornar~\cite{riscv-priv}.

  \item \textbf{Yield a idle para \texttt{move}/\texttt{take}}. El TP3 requería que, tras un \texttt{move} o
  \texttt{take}, la tarea Idle ejecutara el resto del quantum. En x86 eso se implementaba con un \texttt{jmp <idle>}
  explícito. En este port se logra devolviendo \texttt{\&g\_idle\_tf} como próximo \texttt{TrapFrame}: \texttt{trap.S}
  lo cargará y \texttt{sret} dejará a idle corriendo hasta el próximo SSIP.

  \item \textbf{\texttt{getId} sin yield}. También fiel al TP3: retorna al mismo minero con \texttt{tf->a0} seteado.

  \item \textbf{Syscall desconocido}. Se mata a la tarea (\texttt{sched\_kill}), tal como en el TP original.
\end{itemize}

\subsection{Intercambio de tareas por cada tick}

El tick es un trap con \texttt{scause}$=$1 (ver sección~\ref{sec:reloj}). La lógica es:

\begin{codesnippet}
\begin{verbatim}
 if (cause_code == 1) {
     asm volatile("csrc sip, %0" : : "r"((uint64_t)(1ull<<1)));
     if (g_spawn_cooldown_a) g_spawn_cooldown_a--;
     if (g_spawn_cooldown_b) g_spawn_cooldown_b--;
     poll_uart_input();
     update_spinner();
     return schedule_next();
 }
\end{verbatim}
\end{codesnippet}

\texttt{schedule\_next} elige la próxima tarea, \texttt{trap\_handle} retorna el \texttt{TrapFrame*} correspondiente y
\texttt{trap.S} se encarga de restaurarlo y saltar. Todo el ``cambio de tarea'' del TP3 (que en x86 era un \texttt{jmp
<sel>:<off>}) termina siendo una \texttt{sret} con un \texttt{TrapFrame} distinto. No es particularmente elegante, pero
funciona.

\subsection{Atención de excepciones y desalojo}

Cualquier excepción distinta de \texttt{ecall} mata a la tarea actual (si había una corriendo) y cede a idle:

\begin{codesnippet}
\begin{verbatim}
 // Any other exception: kill the current miner (if any) and keep going.
 if (g_current_slot >= 0) {
     print(kHeaderExcPrefix, kHeaderExcX, kHeaderY,
           C_BG_BLACK | C_FG_LIGHT_RED);
     print_hex((uint32_t)cause_code, 2,
               kHeaderExcHexX, kHeaderY,
               C_BG_BLACK | C_FG_LIGHT_RED);
     sched_kill();
 }
 return yield_to_idle();
\end{verbatim}
\end{codesnippet}

Esto cubre page faults (\texttt{scause} 12/13/15), illegal instructions (\texttt{scause}$=$2), misaligned loads/stores,
etc. Por diseño, \texttt{scause} se imprime en hex en el banner superior: es el análogo reducido del ``modo debug'' del
TP3.

\subsection{Modo debug}

El TP3 tenía un modo debug activable con la tecla \texttt{y} que mostraba el estado completo del procesador. En este
port, la información se redujo a un indicador mínimo en la fila 0 del ``video'' (\texttt{TASK EXC XX}) con los dos
dígitos hex del \texttt{cause\_code}. El estado completo de registros no se imprime en pantalla; para debug profundo se
usa QEMU (\texttt{-d int,guest\_errors,cpu\_reset,mmu}). Esta diferencia está documentada en la
sección~\ref{sec:limitaciones}.

\newpage
\section{Despacho de tareas y lógica del juego}
La capa de ``juego propiamente dicho'' es, a nivel código, prácticamente idéntica al TP x86, porque ya era agnóstica a
la ISA (punteros a estructuras, enteros, dibujos de pantalla). Los archivos \texttt{riscv/game.cpp} y
\texttt{riscv/sched.cpp} conservan la lógica de \texttt{game\_move}, \texttt{game\_take}, \texttt{game\_get\_id},
\texttt{sched\_spawn} y \texttt{sched\_kill}.

Las dos diferencias relevantes respecto al informe x86:

\begin{itemize}
  \item \textbf{\texttt{game\_move}} ya no hace \texttt{copy\_page} ni \texttt{mmu\_mapPage} manuales: todo eso está
  encapsulado en \texttt{mmu\_task\_move} (sección~\ref{sec:sv39}). El cuerpo queda enfocado exclusivamente en la lógica
  del juego (cálculo del delta, wrap vertical, espejado para el jugador B, chequeo de base, puntaje, redibujo).

  \item \textbf{\texttt{sched\_spawn}} no registra ningún descriptor de TSS en la GDT (no hay). El \emph{caller} en
  \texttt{kernel.cpp::spawn\_miner} se ocupa de:
  \begin{enumerate}
    \item reservar el slot y marcarlo como vivo (responsabilidad de \texttt{sched\_spawn});
    \item resetear el \texttt{TrapFrame} del slot (\texttt{init\_task\_if\_needed}); y
    \item llamar a \texttt{mmu\_task\_init(slot, player, x, y)}, que construye (o actualiza) la page-table y vuelca el
    binario del minero en la celda de spawn.
  \end{enumerate}
\end{itemize}

\begin{codesnippet}
\begin{verbatim}
 static void spawn_miner(uint32_t player_id) {
     int32_t slot = sched_spawn(player_id);
     if (slot < 0) return;
     const uint32_t u = (uint32_t)slot;
     g_tasks[u].initialized = 0;
     init_task_if_needed(u);
     const uint32_t x = sched_miners_state[u].pos_x;
     const uint32_t y = sched_miners_state[u].pos_y;
     mmu_task_init(u, player_id, x, y);
 }
\end{verbatim}
\end{codesnippet}

El \emph{cooldown} de spawn (\texttt{g\_spawn\_cooldown\_a/b}, 6 ticks) es un detalle UX agregado al port (evita
spawnear 20 mineros apretando una tecla sostenida); no tiene contraparte en el TP3 original.

\newpage
\section{Tareas de usuario}
\label{sec:user}
El TP x86 tenía \texttt{taskA.c}/\texttt{taskB.c} + \texttt{idle.asm}, compilados como blobs que el kernel copiaba desde
\texttt{0x10000} al tablero. Este port sigue el mismo espíritu, pero usa un \emph{build separado} para los binarios de
usuario:

\begin{itemize}
  \item Fuentes: \texttt{riscv/user/miner.c}, \texttt{riscv/user/idle.c}, \texttt{riscv/user/entry.S}.
  \item Linker script propio (\texttt{riscv/user/linker.ld}) con \texttt{.}~$=$~\texttt{0x08000000}.
  \item Producen \texttt{build/user\_miner.bin} / \texttt{build/user\_idle.bin} con \texttt{objcopy -O binary}.
  \item Se embeben en el kernel como símbolos \texttt{\_binary\_build\_user\_miner\_bin\_start} / \texttt{\_end} vía
  \texttt{ld -r -b binary}~\cite{gnu-ld}. El kernel los copia a la celda correspondiente en \texttt{write\_cell\_image}
  durante \texttt{mmu\_task\_init}.
\end{itemize}

El \texttt{entry.S} de usuario es trivial:

\begin{codesnippet}
\begin{verbatim}
 _start:
     .option push
     .option norelax
     la gp, __global_pointer$
     .option pop
     call user_main
 1:  j 1b
\end{verbatim}
\end{codesnippet}

El ABI para invocar servicios del sistema está en \texttt{riscv/syscall.h}, usando la convención \texttt{a7}=id +
\texttt{a0}=arg/retorno (estándar \emph{syscall} de Linux/RISC-V~\cite{riscv-psabi}):

\begin{codesnippet}
\begin{verbatim}
 static inline void syscall_move(e_direction_t d) {
     register uint64_t a7 asm("a7") = 177788;
     register uint64_t a0 asm("a0") = (uint64_t)d;
     asm volatile("ecall" : "+r"(a0) : "r"(a7) : "memory");
 }
\end{verbatim}
\end{codesnippet}

El minero (\texttt{miner.c}) contiene dos modos:

\begin{itemize}
  \item \emph{Cosecha básica}: avanzar, tomar cristales, volver a base (análogo a la tarea ejemplo del TP).
  \item \emph{Prueba de excepción}: cuando el minero tiene \texttt{id}$=$5, provoca una escritura en la VA \texttt{0x0}
  para generar un page fault (equivalente a la ``división por cero'' del TP3).
\end{itemize}

\newpage
\section{Limitaciones y diferencias respecto al TP3 original}
\label{sec:limitaciones}
\begin{itemize}
  \item \textbf{Input por polling, no por IRQ.} El UART 16550 puede generar interrupts vía PLIC, pero se optó por
  polling en cada tick para minimizar la cantidad de mecanismos nuevos a documentar. A efectos de UX no hay diferencia
  perceptible con la frecuencia del tick (${\sim}$20~Hz).
  \item \textbf{Display ANSI, no VGA.} No hay acceso a un framebuffer; el ``tablero'' es una emulación por UART. Esto
  implica cambios de render al redimensionar la terminal, y que el archivo de video no es un ``segmento de pantalla''
  \emph{stricto sensu}.
  \item \textbf{Modo debug simplificado.} Como se comentó en el ejercicio~7, la pantalla de debug completa del TP3
  (registros, flags, stack) no se portó.
  \item \textbf{Sin OpenSBI.} Al correr con \texttt{-bios none}, el kernel implementa su propia firmware mínima
  (\texttt{machine.S}). Con \texttt{-bios default} (OpenSBI) el mismo kernel podría reemplazar \texttt{machine.S} por
  llamadas SBI (\texttt{sbi\_set\_timer}, \texttt{sbi\_console\_putchar}, etc.) sin cambios a la lógica en S-mode.
  \item \textbf{Single-hart.} QEMU \texttt{virt} soporta SMP, pero este port corre sobre un solo HART. Extenderlo a SMP
  requeriría IPI (via \texttt{msip}), serialización del scheduler, y probablemente per-hart scratch state.
  \item \textbf{Mapa de memoria distinto al del enunciado.} La figura del enunciado asume RAM desde \texttt{0x0000}. En
  QEMU \texttt{virt} la RAM arranca en \texttt{0x8000\_0000}, así que las direcciones físicas concretas (kernel,
  page-tables, tablero) son otras; lo que se preserva es la \emph{estructura} del layout.
  \item \textbf{3 niveles de page-table vs 2.} Sv39 tiene L2/L1/L0; el TP3 x86 (paginación 32 bits sin PAE) tenía PD/PT.
  Esto se aprovecha para mapear RAM con una sola gigapage, algo imposible en IA-32 de 32 bits sin PAE.
\end{itemize}

\newpage
\section{Build y ejecución}
\subsection{Dependencias}

\begin{itemize}
  \item \texttt{make}.
  \item Toolchain bare-metal: \texttt{riscv64-unknown-elf-\{gcc, g++, as, ld, objcopy\}}.
  \item \texttt{qemu-system-riscv64}.
\end{itemize}

\subsection{Compilación}

\begin{codesnippet}
\begin{verbatim}
 make -C riscv clean all
\end{verbatim}
\end{codesnippet}

Genera \texttt{riscv/build/kernel.elf} (y \texttt{.bin}), que contiene embebidos los binarios de usuario
(\texttt{build/user\_miner.bin}, \texttt{build/user\_idle.bin}).

\subsection{Ejecución}

\begin{codesnippet}
\begin{verbatim}
 make -C riscv run
\end{verbatim}
\end{codesnippet}

que se expande a:

\begin{codesnippet}
\begin{verbatim}
 qemu-system-riscv64 -machine virt -nographic -monitor none \
     -serial stdio -bios none -kernel riscv/build/kernel.elf
\end{verbatim}
\end{codesnippet}

\subsection{Controles}

\begin{itemize}
  \item Jugador ROJO: \texttt{W} (arriba), \texttt{S} (abajo), \texttt{Q} o \texttt{E} (spawnear minero).
  \item Jugador AZUL: \texttt{I} (arriba), \texttt{K} (abajo), \texttt{O} o \texttt{P} (spawnear minero).
\end{itemize}

En la esquina superior derecha se muestra, en hex, el último byte recibido por el UART (útil para validar input).

\newpage

% ============================================================
\begin{thebibliography}{99}

\bibitem{riscv-unpriv}
RISC-V International. \emph{The RISC-V Instruction Set Manual, Volume I: Unprivileged Architecture}. \\
\url{https://riscv.github.io/riscv-isa-manual/snapshot/unprivileged/}

\bibitem{riscv-priv}
RISC-V International. \emph{The RISC-V Instruction Set Manual, Volume II: Privileged Architecture}. Document version
20211203, ratificado en diciembre de 2021. \\
\url{https://github.com/riscv/riscv-isa-manual/releases/download/Priv-v1.12/riscv-privileged-20211203.pdf} \\
Versión snapshot en HTML: \url{https://riscv.github.io/riscv-isa-manual/snapshot/privileged}

\bibitem{riscv-psabi}
RISC-V Non-ISA Working Group. \emph{RISC-V ELF psABI Specification}. \\
\url{https://riscv-non-isa.github.io/riscv-elf-psabi-doc/}

\bibitem{riscv-sbi}
RISC-V International. \emph{RISC-V Supervisor Binary Interface Specification}. \\
\url{https://github.com/riscv-non-isa/riscv-sbi-doc}

\bibitem{riscv-aclint}
RISC-V International. \emph{RISC-V Advanced Core Local Interruptor (ACLINT) Specification}. \\
\url{https://github.com/riscv/riscv-aclint}

\bibitem{riscv-plic}
RISC-V International. \emph{RISC-V Platform-Level Interrupt Controller (PLIC) Specification}. \\
\url{https://github.com/riscv/riscv-plic-spec}

\bibitem{qemu-virt}
QEMU Project. \emph{QEMU System Emulation, 'virt' Machine (RISC-V)}. \\
\url{https://www.qemu.org/docs/master/system/riscv/virt.html}

\bibitem{xv6-riscv}
Cox, R., Kaashoek, M.~F., Morris, R. \emph{xv6: a simple, Unix-like teaching operating system (RISC-V port)}. MIT. \\
\url{https://github.com/mit-pdos/xv6-riscv} \\
\url{https://github.com/mit-pdos/xv6-riscv-book}

\bibitem{gnu-ld}
Free Software Foundation. \emph{Using ld (GNU linker) --- Binary input format}. \\
\url{https://sourceware.org/binutils/docs/ld/Options.html}

\bibitem{ns16550}
National Semiconductor. \emph{PC16550D Universal Asynchronous Receiver/Transmitter with FIFOs}. Datasheet.

\bibitem{osdev-riscv}
OSDev Wiki. \emph{RISC-V} (página principal del wiki), \emph{RISC-V Bare Bones}, y \emph{RISC-V Meaty Skeleton with QEMU
virt board}. \\
\url{https://wiki.osdev.org/RISC-V} \\
\url{https://wiki.osdev.org/RISC-V_Bare_Bones} \\
\url{https://wiki.osdev.org/RISC-V_Meaty_Skeleton_with_QEMU_virt_board}

\end{thebibliography}

\end{document}
